% EDITING: type inside the final {...} of each field. Change the shown font size
% (for example, font=10pt) in that same field block when needed.

\vspace{1.1cm}

\begingroup
\renewcommand{\clearpage}{}
\GrantSection{3}{Research Strategy / Pioneer-Style Essay (Next Page)}
\endgroup
\GrantWarning{RFA-RM-27-001 uses a five-page Research Strategy essay and instructs applicants not to submit Specific Aims or a detailed scientific plan. For another funder, rename, reorder, add, or delete the page modules to match that sponsor's instructions.}

% \GrantSection{3}{Research Strategy / Pioneer-Style Essay}
% \GrantWarning{RFA-RM-27-001 uses a five-page Research Strategy essay and instructs applicants not to submit Specific Aims or a detailed scientific plan. For another funder, rename, reorder, add, or delete the page modules to match that sponsor's instructions.}

\GrantResearchPage{1}{
  \GrantGuidance{Begin with the two scientific-area designations and a descriptive project title. Explain the major challenge or opportunity, importance, and broad scientific vision for a multidisciplinary audience.}
  \GrantLineField[id=essay-science-areas,font=10pt,maxchars=80]{Scientific Areas (primary; secondary)}{3 CTR, 5 IE}{\DefaultScienceAreas}
  \GrantLineField[id=essay-project-title,font=11pt,maxchars=220]{Project Title}{Daraxonrasib Phase 1 LLM-Directed Robotic Whipple in KRAS-Mutated PDAC}
  \GrantTextArea[
  id=essay-project-description-1, font=8pt, height=5.7in, maxchars=5000
]{Project description - major challenge or opportunity and why it matters}{Pancreatic ductal adenocarcinoma (PDAC) presents a convergence of unmet biological, surgical, and translational challenges. Even when disease is technically resectable, curative-intent pancreaticoduodenectomy is complex, perioperative morbidity can compromise treatment, and recurrence remains common. The proposed project addresses this problem in patients with resectable ^^Jor borderline-resectable KRAS G12-mutated PDAC by integrating perioperative daraxonrasib (RMC-6236), an oral RAS(ON) multi-selective inhibitor, with an on-premises large-language-model-directed eight-arm robotic Whipple procedure \citep{phase1ind,onpremwhippl,daraxwhipple}. ^^J ^^JThe major opportunity is not simply to place an artificial-intelligence layer on an existing operation. It is to create a new clinical-trial architecture in which molecular therapy, robotic execution, real-time sensing, human supervision, regulatory documentation, and computational assurance are designed as one auditable system. Current oncology development often treats these domains as separate handoffs. Drug development proceeds through one regulatory spine, devices through another, clinical operations through site-specific systems, and software verification through disconnected technical reviews. The proposed work instead uses ^^Jharmonized IND and IDE logic, ICH E6(R3)-aligned trial conduct, and an explicit Physical AI assurance layer so that the more stringent applicable requirement governs each patient-facing activity \citep{onpremwhippl,cfr312paiuni,cfr050paiuni,ichpaistduni}. ^^J ^^JThis matters because the central barrier to responsible clinical deployment is not whether an LLM can generate code, documents, or robotic commands. The barrier is whether every consequential output can be verified before it reaches a participant. The investigator's prior work therefore places verification, validation, and uncertainty quantification (VVUQ) ahead of generation. In that framework, a proposed action is accepted only when all required checks pass; otherwise it is escalated to a qualified human or blocked. Repository-scale demonstrations have used deterministic test suites, ten-gate assurance structures, reproducible seeds, audit records, and explicit ACCEPT, ESCALATE, or BLOCK outcomes \citep{vvuqllmverif,clintrustpai,h2pancrevvuq}. ^^JThese studies are not evidence of clinical readiness, but they establish a testable engineering premise for a governed first-in-human program. ^^J ^^JThe opportunity is also operational. A Phase 1 oncology program depends on a large sequence of interlocking documents, including the IND and IRB package, protocol and consent revisions, cohort-review packages, clinical-hold responses, and later regulatory modules. Repository-based LLM workflows can organize those materials through traceable draft-to-final pipelines while preserving human review and fixed regulatory clocks \citep{phase1guidance,oncotrialllm,oncotrialeff}. This administrative acceleration could reduce avoidable preparation delays without shortening participant follow-up, bypassing regulatory review, or substituting generated text for accountable scientific judgment. ^^J ^^JThe proposed research therefore asks a high-value question: can a combined drug, robotic, and computational-assurance platform be translated from auditable simulation and regulatory-document prototypes into a carefully bounded Phase 1 study that establishes safety, feasibility, and readiness for later comparative testing? Success would create a generalizable pathway for evaluating other Physical AI oncology interventions; failure would still yield valuable stopping criteria, assurance evidence, and regulatory lessons. Either outcome would clarify how advanced robotics and AI can be introduced into cancer trials without lowering the standards ^^Jowed to participants.
}
}


\GrantResearchPage{2}{
  \GrantGuidance{Describe the pioneering approach and potential for groundbreaking or paradigm-shifting results. Convey rigor and robustness without presenting a conventional step-by-step protocol. Essential citations and figures, if used, consume space within the five-page limit.}
  \GrantTextArea[
  id=essay-project-description-2, font=8.35pt, height=6.9in, maxchars=6200
]{Project description - innovative vision, conceptual framework, rigor, and potential impact}{The innovative vision is a verification-first clinical-trial system in which artificial intelligence is never treated as an autonomous source of truth. The conceptual framework has five linked layers: patient and disease characterization; perioperative daraxonrasib treatment; robotic surgical execution; human clinical authority; and an auditable trust infrastructure that ^^Jevaluates every software-generated or robot-executed action. The project is innovative because these layers are co-designed rather than appended after the fact. The Phase 1 application defines the clinical population and combined intervention, while the companion protocol frames simultaneous drug, device, and Good Clinical Practice obligations \citep{phase1ind,onpremwhippl}. ^^J ^^JAt the center of the framework is verification before generation. A software agent may propose an action, but it cannot authorize itself. Predefined VVUQ gates evaluate computational correctness, boundary conditions, sensor consistency, uncertainty, reproducibility, and conformance with the approved protocol. The disposition is explicit: ACCEPT when all ^^Jrequired criteria are met, ESCALATE when qualified human review is required, or BLOCK when safety or validity criteria fail \citep{clintrustpai,vvuqllmverif,h2pancrevvuq}. Human investigators retain legal accountability, procedure authority, emergency-stop control, and the ability to override or terminate automation. This architecture is designed to counter both automation bias and algorithm aversion by making reliance proportional to demonstrated capability. ^^J ^^JRigor will be established through a prospective assurance case rather than retrospective narrative. Before clinical activation, the project will freeze the intended context of use, software and model versions, sensor interfaces, command boundaries, prohibited actions, human-authorization points, and failure-response pathways. Verification will include unit, integration, boundary, timing, emergency-stop, fail-safe, and deterministic-replay tests. Validation will compare system behavior with approved surgical and clinical requirements. Uncertainty quantification will characterize variability in sensing, timing, model outputs, and simulated patient states. Provenance records will link each input, model version, generated artifact, review, decision, and correction \citep{natlmcpservr,fedlearnpaio,paistdlevelu}. ^^J ^^JClinical rigor will be equally explicit. The initial study will remain first-in-human, prospective, open-label, and single-arm. Its purpose is to estimate safety, tolerability, procedural feasibility, and dose or procedure readiness, not to make definitive ^^Jefficacy claims \citep{phase1ind,daraxwhipple}. Prespecified analysis populations, missing-data rules, adjudication procedures, stopping boundaries, and sensitivity analyses will be finalized before database lock. A medical monitor and independent data and safety monitoring body will review cumulative drug, surgical, and cyber-physical risk. ^^J ^^JThe platform will also test a new method for trial execution. Repository-based LLMs can generate synchronized protocol, consent, regulatory, and operational drafts through a traceable multi-file workflow, but every output will remain subject to source verification, version control, role-based approval, and formal change management \citep{phase1guidance,oncotrialllm,paiautospons}. Standardized authorization, FHIR clinical-data, DICOM imaging, audit-ledger, and provenance services ^^Jcan reduce fragmentation while maintaining privacy and site control \citep{natlmcpservr,fedlearnpaio}. These tools are supporting infrastructure, not substitutes for investigators, IRBs, regulators, statisticians, surgeons, pharmacists, participants.^^J^^JThe potential impact is broad. A successful Phase 1 program would establish an evidence-based route for combining targeted therapy with governed Physical AI in surgery, define practical regulatory and assurance requirements, and create reusable templates for future oncology trials. It could also show that speed and rigor are not opposing goals: document preparation, simulation, and quality checks may be accelerated while clinical follow-up, independent review, and patient protections remain intact \citep{phase1guidance,oncotrialeff}. A negative or mixed result would be equally informative by identifying unacceptable failure modes, infeasible interfaces, or insufficient assurance margins before wider deployment. The intended paradigm shift is therefore from AI-enabled generation to accountable, verification-governed clinical action.
}
}

\GrantResearchPage{3}{
  \GrantGuidance{Provide concrete evidence of the PD/PI's innovativeness, such as integrating disparate information, questioning paradigms, working across disciplines, handling uncertainty, persisting through failure, and executing ambitious strategies.}
  \GrantTextArea[
  id=essay-innovativeness, font=9pt, height=5.63in, maxchars=6200
]{Evidence of PD/PI innovativeness}{Kevin Kawchak's recent work demonstrates a sustained pattern of integrating disciplines that are usually separated: oncology drug development, pancreatic surgery, regulatory writing, robotic systems, software assurance, quantitative systems pharmacology, digital twins, federated learning, and repository-scale language models. The portfolio moves from literature-grounded cancer analyses and in silico PDAC studies to protocol, IND, site, and sponsor-system concepts, creating a connected body of work rather than isolated demonstrations \citep{pdacdigtwinp,chatgpt100kp,qspmetpancre,phase1ind,onpremwhippl}. ^^J ^^JA central example is the shift from generating outputs to interrogating whether those outputs are trustworthy. The VVUQ work explicitly argues that code generation is not the decision-bearing accomplishment; verification, validation, uncertainty quantification, and reproducible execution are. Demonstrations include multi-module automated testing, deterministic replay, numerical sensitivity analyses, and decision gates that accept, escalate, or block proposed actions \citep{vvuqllmverif,h2pancrevvuq,fdadigtwinpc}. This is a creative reframing of the AI problem from productivity to accountable evidence. ^^J ^^JThe investigator has also pursued unusually broad implementation artifacts: a combined Phase 1 IND and protocol concept; a Phase 2 comparative protocol; repository-based document-generation guidance; an AI-native sponsor architecture with retained human accountability; standardized MCP services for authorization, clinical data, imaging, audit, and provenance; and a federated-learning framework designed around privacy and cross-site cooperation \citep{phase1ind,daraxwhipple,phase1guidance,paiautospons,natlmcpservr,fedlearnpaio}. These works show persistence in converting ideas into auditable documents, code, tests, diagrams, and operational frameworks. ^^J ^^JThe record also shows willingness to work under uncertainty. Simulation studies explicitly identify calibration limits, synthetic-patient gaps, hypothetical-hardware gaps, and the need for prospective validation rather than presenting computational results as clinical evidence \citep{qspmetpancre,fdadigtwinpc,pdac060s2030}. The proposed project extends that discipline by defining clinical translation as a falsifiable test with stopping rules, independent oversight, and transparent failure reporting. ^^J ^^JThis combination of cross-disciplinary synthesis, verification-first reasoning, rapid but traceable execution, and explicit treatment of limitations is directly relevant to a Pioneer-style award. The investigator's distinctive ^^Jcontribution is not a claim that one model or robot is already ready for patient care. It is the construction of an integrated research direction in which advanced generation tools are subordinated to clinical governance, evidence, and participant protection.
}
}

\GrantResearchPage{4}{
  \GrantGuidance{Explain why the proposed work is a fundamental new direction rather than an obvious extension or scale-up of current or past work. Then explain why a high-risk/high-reward or otherwise special mechanism is a better fit than a conventional award.}
  \GrantTextArea[
  id=essay-new-direction, font=10pt, height=4.2in, maxchars=3400
]{How the planned research substantially differs from past or current work}{The planned research is a fundamental new direction because the investigator's prior portfolio is predominantly computational, documentary, regulatory-conceptual, or simulation-based. It includes ^^JPDAC digital-twin proposals, large in silico trials, QSP models, VVUQ test frameworks, robotic simulations, protocol drafts, and AI-assisted sponsor and site architectures \citep{pdacdigtwinp,chatgpt100kp,qspmetpancre,vvuqllmverif,pdac060s2030,paiautospons}. Those works establish hypotheses, workflows, and assurance methods, but don't establish safety or feasibility in participants. ^^J ^^JThe proposed project crosses that boundary. It converts a body of simulated and document-based work into a prospective, first-in-human clinical investigation of perioperative daraxonrasib with an LLM-directed robotic Whipple in resectable or borderline-resectable KRAS G12-mutated PDAC \citep{phase1ind,onpremwhippl}. This requires real regulatory authorization, qualified clinical sites, investigational ^^Jproduct control, device accountability, surgeon and staff credentialing, informed consent, independent monitoring, validated data systems, and direct responsibility for participant welfare. The scientific ^^Jquestion is no longer whether an AI system can generate a plausible protocol or complete a synthetic test suite. It is whether a bounded, human-governed version of the combined platform can meet prespecified clinical and cyber-physical safety requirements. ^^J ^^JThe work also differs from a conventional scale-up. It does not merely add patients, compute, prompts, ^^Jor robotic arms to a prior simulation. The clinical context of use, risk classification, assurance case, endpoints, stopping boundaries, and human-machine authority structure must be redesigned for prospective use. The transition from simulated performance to accountable human-subject research is therefore the core innovation and the principal source of risk.
}
  \GrantTextArea[
  id=essay-program-fit, font=10pt, height=2.85in, maxchars=2300
]{Suitability for the target program / why a traditional mechanism is insufficient}{This project fits a high-risk, high-reward Pioneer mechanism because it proposes a potentially field changing clinical-trial architecture rather than an incremental optimization of an established procedure. The concept combines targeted KRAS therapy, complex pancreatic surgery, Physical AI, human ^^Joversight, regulatory harmonization, and VVUQ into a single prospective program \citep{phase1ind,onpremwhippl,clintrustpai}. The scientific and operational uncertainties are too interdependent ^^Jfor a narrowly scoped conventional award: the value of the drug-device concept cannot be separated from the assurance framework, site infrastructure, regulatory pathway, and real-time governance needed. ^^J ^^JA traditional mechanism would also encourage fragmentation into sequential subprojects, increasing the risk that the clinical protocol, robotic system, data architecture, and safety case evolve incompatibly. The Pioneer structure supports a unified program with authority to stop, redesign, or redirect components as evidence emerges. It is especially appropriate because the most important deliverable may be a ^^Jvalidated translational framework and a clear boundary on what should not proceed, rather than a predetermined positive efficacy result. The proposal aligns with the program's emphasis on unusually innovative research, broad impact, and a substantial change in direction \citep{nih_rfa_rm_27_001}.
}
}

\GrantResearchPage{5}{
  \GrantGuidance{For the NIH Pioneer baseline, include the required effort commitment: more than six person-months (at least 51 percent) in years 1-3, at least four months (33 percent) in year 4, and at least three months (25 percent) in year 5. Explain adjustment of existing commitments when applicable.}
  \GrantTextArea[
  id=essay-effort-commitment, font=10pt, height=4.0in, maxchars=2600
]{Statement of research effort commitment and adjustment of existing commitments}{If selected for award, Kevin Kawchak will commit more than 6 person-months of research effort, equivalent to at least 51 percent, in each of project years 1 through 3; at least 4 person-months, 33 percent, in project year 4; and at least 3 person-months, 25 percent, in project year 5, consistent with the Pioneer Award requirements \citep{nih_rfa_rm_27_001}. To make that commitment operationally credible, current non-sponsored and internally directed activities at ChemicalQDevice will be consolidated into the funded project, while routine website, general portfolio-maintenance, and nonessential exploratory drafting work will be reduced, deferred, or delegated. During the first three years, calendar priority will be assigned to regulatory filings, site-startup oversight, safety review, documentation governance, statistical review, and manuscript-ready curation of the trial platform; external talks, speculative side projects, and unrelated consulting will be capped or paused.^^J^^JRepository-based workflows, versioned templates, and re-use of already established IND/protocol/ simulation assets will further reduce avoidable administrative overhead while preserving accountable human review \citep{phase1guidance,oncotrialllm,oncotrialeff,phase1ind,onpremwhippl}. Based on the materials provided for this application draft, no conflicting active federal award requiring 6 ^^Jor more person-months was identified; if any additional current or pending support exists at time of submission, the final other-support statement should be revised so that the exact source, overlap ^^Janalysis, rebudgeting, delegation, no-cost extension, or relinquishment steps are stated explicitly. ChemicalQDevice will provide institutional confirmation that protected effort, essential decision authority, secure compute resources, and records stewardship needed for the project are available across the full project period.
}
  \GrantTextArea[
  id=essay-conclusion, font=10pt, height=4.45in, maxchars=3300
]{Concluding case for broad impact, feasibility, responsible risk management, and measurable success}{This project seeks to establish whether a verification-governed combination of perioperative daraxonrasib and an LLM-directed robotic Whipple can enter early clinical testing without weakening the protections expected in oncology research. Its broad impact lies in the framework: a combined drug device protocol, explicit human authority, harmonized regulatory responsibilities, auditable provenance, and VVUQ gates that accept, escalate, or block patient-facing actions \citep{phase1ind,onpremwhippl,clintrustpai,vvuqllmverif}. ^^J ^^JFeasibility is supported by a coherent body of preparatory work spanning PDAC simulation, QSP and digital-twin methods, Phase 1 and Phase 2 protocol drafting, repository-based document workflows, federated data infrastructure, and standardized authorization and provenance services \citep{qspmetpancre,fdadigtwinpc,phase1guidance,daraxwhipple,fedlearnpaio,natlmcpservr}. These ^^Jproducts do not prove clinical performance; they reduce ambiguity about the questions, interfaces, documents, tests, and governance that must be resolved before enrollment. Risk will be managed ^^Jthrough staged decision gates. During conduct, prespecified stopping rules will address drug toxicity, surgical complications, device failures, cyber-physical anomalies, protocol deviations, and unacceptable uncertainty. Human clinicians retain authority at every stage. ^^J ^^JMeasurable success will be defined by regulatory and site activation; completion of prespecified verification and validation tests; enrollment and retention against approved targets; interpretable ^^Jestimates of safety and feasibility; complete auditability of patient-facing automated actions; timely safety review; and a final decision to proceed, redesign, or stop. The project is deliberately structured so that an unfavorable result remains scientifically valuable. By replacing unbounded enthusiasm for AI generation with prospective evidence and accountable control, the work could establish a durable standard for responsible Physical AI oncology trials.
}
}
